\documentclass[aspectratio=169,11pt]{beamer}
\usepackage{booktabs}
\usepackage{pgfplots}
\usepackage{pgfplotstable}
\usepackage{amsmath}
\usepackage{array}

\pgfplotsset{compat=1.18}
% ============================================================
% Synthetic Library Space Utilisation Dataset
% ============================================================
%
% This dataset is fully synthetic and was constructed specifically
% for this benchmark asset. It is NOT presented as real university data.
%
% Observation unit:
% One study zone x one time period x one day.
%
% Observation period:
% One hypothetical week, Monday to Saturday.
%
% Study zones and capacities:
% Quiet Study Zone   = 40 seats
% General Study Zone = 60 seats
% Group Study Zone   = 30 seats
%
% Time periods:
% P1 = 8:00-10:00 AM
% P2 = 10:00 AM-12:00 PM
% P3 = 2:00-4:00 PM
% P4 = 4:00-6:00 PM
%
% Occupancy rate is derived as:
% OccupiedSeats / Capacity * 100
% ============================================================

\pgfplotstableread[col sep=comma]{
Day,Period,Zone,Capacity,OccupiedSeats,AvgStudyDuration,StudentEntries,SatisfactionScore
Monday,P1,Quiet,40,28,85,24,3.3
Monday,P1,General,60,38,93,34,3.5
Monday,P1,Group,30,16,79,9,3.3
Monday,P2,Quiet,40,31,92,17,3.1
Monday,P2,General,60,46,51,30,3.0
Monday,P2,Group,30,19,74,14,3.3
Monday,P3,Quiet,40,28,78,18,3.2
Monday,P3,General,60,44,102,36,3.4
Monday,P3,Group,30,17,102,9,3.8
Monday,P4,Quiet,40,37,80,19,3.0
Monday,P4,General,60,50,101,33,3.3
Monday,P4,Group,30,21,53,10,2.3
Tuesday,P1,Quiet,40,29,89,27,2.9
Tuesday,P1,General,60,41,96,30,3.3
Tuesday,P1,Group,30,16,80,15,3.4
Tuesday,P2,Quiet,40,32,69,27,2.9
Tuesday,P2,General,60,49,61,34,3.0
Tuesday,P2,Group,30,21,106,14,3.8
Tuesday,P3,Quiet,40,31,62,21,2.5
Tuesday,P3,General,60,42,81,27,3.2
Tuesday,P3,Group,30,18,47,9,2.2
Tuesday,P4,Quiet,40,36,52,26,3.0
Tuesday,P4,General,60,55,99,45,3.6
Tuesday,P4,Group,30,22,49,12,2.9
Wednesday,P1,Quiet,40,29,64,23,2.5
Wednesday,P1,General,60,41,58,32,3.1
Wednesday,P1,Group,30,17,80,16,3.0
Wednesday,P2,Quiet,40,34,57,17,2.6
Wednesday,P2,General,60,48,47,36,2.6
Wednesday,P2,Group,30,19,103,11,3.6
Wednesday,P3,Quiet,40,31,76,21,3.1
Wednesday,P3,General,60,46,49,28,3.0
Wednesday,P3,Group,30,17,53,17,3.1
Wednesday,P4,Quiet,40,37,91,24,3.6
Wednesday,P4,General,60,56,72,56,3.0
Wednesday,P4,Group,30,23,99,21,3.1
Thursday,P1,Quiet,40,33,79,22,3.4
Thursday,P1,General,60,36,58,36,3.1
Thursday,P1,Group,30,18,54,8,3.1
Thursday,P2,Quiet,40,36,83,27,3.3
Thursday,P2,General,60,47,76,28,2.9
Thursday,P2,Group,30,19,75,14,2.9
Thursday,P3,Quiet,40,30,74,20,3.4
Thursday,P3,General,60,41,74,27,2.8
Thursday,P3,Group,30,17,77,9,3.4
Thursday,P4,Quiet,40,38,60,26,3.1
Thursday,P4,General,60,52,75,27,3.5
Thursday,P4,Group,30,24,73,17,3.5
Friday,P1,Quiet,40,29,72,17,3.4
Friday,P1,General,60,37,67,31,3.3
Friday,P1,Group,30,15,59,15,3.0
Friday,P2,Quiet,40,33,67,16,2.7
Friday,P2,General,60,44,46,45,2.5
Friday,P2,Group,30,20,101,13,3.5
Friday,P3,Quiet,40,28,99,26,3.4
Friday,P3,General,60,37,100,28,3.1
Friday,P3,Group,30,18,61,8,2.6
Friday,P4,Quiet,40,36,106,17,3.7
Friday,P4,General,60,50,107,47,3.4
Friday,P4,Group,30,22,91,16,3.3
Saturday,P1,Quiet,40,27,72,22,3.1
Saturday,P1,General,60,30,48,14,2.7
Saturday,P1,Group,30,12,107,8,4.0
Saturday,P2,Quiet,40,32,92,30,3.7
Saturday,P2,General,60,39,56,18,2.6
Saturday,P2,Group,30,18,75,8,3.0
Saturday,P3,Quiet,40,27,89,28,3.1
Saturday,P3,General,60,40,106,34,3.3
Saturday,P3,Group,30,15,76,16,3.3
Saturday,P4,Quiet,40,33,72,23,3.1
Saturday,P4,General,60,50,108,29,3.3
Saturday,P4,Group,30,19,45,15,2.9
}\librarydata
\usetheme{Madrid}

\title{Analysing University Library Space Utilisation and Student Study Patterns}
\author{Oluwadurotimi Alade}
\institute{University of Ibadan}
\date{September 2026}

\begin{document}

\begin{frame}
    \titlepage
\end{frame}

\begin{frame}{Presentation Overview}
    \tableofcontents
\end{frame}

\section{Introduction}

\begin{frame}{Background}
    \begin{itemize}
        \item University libraries provide spaces for individual and collaborative study.
        \item The way students use these spaces can vary across zones, time periods, and days.
        \item Understanding these patterns can support better use of available study spaces.
    \end{itemize}
\end{frame}

\begin{frame}{Problem Statement}
    Library study spaces have limited seating capacity, but demand may vary
    across different zones, time periods, and days. A structured analysis of
    occupancy, study duration, student entries, and satisfaction can help
    identify patterns in space utilisation.
\end{frame}

\begin{frame}{Aim and Objectives}
    \textbf{Aim}

    To analyse university library space utilisation and student study patterns
    using synthetic observational data.

    \vspace{0.4cm}

    \textbf{Objectives}
    \begin{enumerate}
        \item Measure occupancy rates across library study zones.
        \item Compare utilisation across time periods and days.
        \item Examine study duration, student entries, and satisfaction.
        \item Explore relationships among selected study-pattern variables.
        \item Develop practical recommendations from the observed patterns.
    \end{enumerate}
\end{frame}

\section{Research Questions}

\begin{frame}{Research Questions}
    \begin{enumerate}
        \item How does library space occupancy vary across the three study zones?
        \item How does occupancy vary across different time periods?
        \item How does occupancy vary across days of the week?
        \item What patterns exist between student entries, study duration, and occupancy?
        \item What relationship exists between occupancy and student satisfaction?
    \end{enumerate}
\end{frame}

\section{Data and Methodology}

\begin{frame}{Study Design}
    \begin{itemize}
        \item Data type: fully synthetic benchmark data.
        \item Observation period: one hypothetical week, Monday to Saturday.
        \item Study zones: Quiet, General, and Group.
        \item Four two-hour observation periods were used each day.
        \item Total observations: 72.
    \end{itemize}
\end{frame}

\begin{frame}{Variables}
    \begin{table}
        \centering
        \small
        \begin{tabular}{ll}
            \toprule
            \textbf{Variable} & \textbf{Description} \\
            \midrule
            Day & Day of observation \\
            Period & Two-hour observation period \\
            Zone & Library study zone \\
            Capacity & Available seats in the zone \\
            Occupied Seats & Seats occupied during observation \\
            Study Duration & Average study duration in minutes \\
            Student Entries & Student entries recorded during the period \\
            Satisfaction & Average satisfaction score from 1 to 5 \\
            \bottomrule
        \end{tabular}
    \end{table}
\end{frame}

\begin{frame}{Observation Periods}
    \begin{table}
        \centering
        \small
        \begin{tabular}{cl}
            \toprule
            \textbf{Period} & \textbf{Time} \\
            \midrule
            P1 & 8:00--10:00 AM \\
            P2 & 10:00 AM--12:00 PM \\
            P3 & 2:00--4:00 PM \\
            P4 & 4:00--6:00 PM \\
            \bottomrule
        \end{tabular}
    \end{table}

    \vspace{0.3cm}

    Each observation covers one study zone during one two-hour period
    on one day.
\end{frame}

\section{Preliminary Results}

\begin{frame}{Preliminary Occupancy Result}
    The current validated dataset gives an overall mean occupancy rate of:

    \[
    \boxed{42.60\%}
    \]
\end{frame}

\begin{frame}
    \centering
    \Huge Questions?
\end{frame}

\end{document}
